\section{Derivations connecting the projection audit to overlap preservation}
\label{qdo:appendix}

This appendix separates three objects: the historical training intervention,
the information in a fixed linear measurement, and the information remaining
after that measurement is rectified. The Gaussian calculation supplies an
exact diagnostic for the original synthetic task. Its application to a
selective overlap channel requires an additional conditioning argument.

\subsection{Historical capacity, information loss, and anchor geometry}
\label{qdo:historical}

The archived student was a bias-free network
\begin{equation}
 f_\theta(x)=\sum_{i=1}^{H}a_i\ReLU(w_i^\top x),
 \qquad W=[w_1,\ldots,w_H]\in\R^{8\times H}.
 \label{qdo:student}
\end{equation}
Its teacher used four orthonormal directions:
\begin{equation}
 f_*(x)=\sum_{j=1}^{4}c_j\ReLU(q_j^\top x),\qquad
 c=(-\tfrac12,+\tfrac12,-\tfrac12,+\tfrac12),\qquad
 q_i^\top q_j=\delta_{ij}.
 \label{qdo:teacher}
\end{equation}
Thus width four had enough capacity to express the scalar target. Nevertheless,
every $8$-to-$4$ linear map satisfies
\begin{equation}
 \dim\ker W^\top=8-\rank W\ge4.
 \label{qdo:nullity}
\end{equation}
For any $v\in\ker W^\top$, the inputs $x$ and $x+v$ have identical
preactivations and identical post-ReLU states. Preserving every input was
therefore impossible in this architecture. Preserving the target was possible
because it depends only on the four-dimensional teacher span. A successful
representation could also discard distinctions within that span whenever they
leave the target unchanged. The audit concerned collisions with \emph{different}
required targets, rather than every collision or an ordinary rank reduction
\cite{qdo,projection}.

The historical QDO intervention acted on parameters. On the domain
$r_i=\norm{w_i}>0$, define
\begin{equation}
 u_i=w_i/r_i,\qquad b_i=a_i r_i,\qquad
 \phi(\theta)=\operatorname{pack}(u_1,\ldots,u_H,b_1,\ldots,b_H).
 \label{qdo:canonical}
\end{equation}
For positive $s_i$, replacing $(w_i,a_i)$ by $(s_iw_i,a_i/s_i)$ leaves both
$f_\theta$ and $\phi(\theta)$ unchanged, by positive homogeneity of ReLU.
The coordinates intentionally remove this parameter ambiguity; they are not
an encoding from which every original parameter value can be reconstructed.

With fixed historical parameters $\alpha_j$, an orthonormal-column matrix
$Q\in\R^{N\times k}$, and bandwidth $\tau>0$, write
\begin{align}
 z&=Q^\top\phi(\theta),& z_j&=Q^\top\phi(\alpha_j),&
 \delta_j&=z-z_j,\\
 \mathcal R(\theta)&=\sum_j\omega_j,&
 \omega_j&=\exp\!\left(-\frac{\norm{\delta_j}^2}{2\tau^2}\right),&
 r_z&=\frac{1}{\tau^2}\sum_j\omega_j\delta_j.
 \label{qdo:anchor}
\end{align}
The negative gradient is
\begin{equation}
 F=-\nabla_\theta\mathcal R=D\phi(\theta)^\top Qr_z.
 \label{qdo:force}
\end{equation}
For $v=Qr_z$, unpacked into $v_{u_i},v_{b_i}$, differentiation gives
\begin{align}
 \mathrm du_i&=r_i^{-1}(I-u_i u_i^\top)\,\mathrm dw_i,&
 \mathrm db_i&=a_i u_i^\top\mathrm dw_i+r_i\,\mathrm da_i,\\
 F_{w_i}&=r_i^{-1}(I-u_i u_i^\top)v_{u_i}+a_i v_{b_i}u_i,&
 F_{a_i}&=r_i v_{b_i}.
 \label{qdo:pullback}
\end{align}
Consequently $w_i^\top F_{w_i}-a_iF_{a_i}=0$: this auxiliary force is
orthogonal to the infinitesimal positive-rescaling direction. Penalty
invariance does not establish invariance of the complete Adam trajectory,
nor does it establish injectivity of a forward representation \cite{qdo}.
The optimizer projection $Q^\top\phi(\theta)$ and input measurement $W^\top x$
act on different spaces. The proposal leaves this optimizer intervention
behind and places its preservation requirement on the forward state and its
overlap records.

This distinction also separates their possible guarantees. An auxiliary
parameter force can change which representation is learned; it does not give
a decoder access to information missing from a particular frozen forward
state. An overlap record instead changes what that decoder receives. Its
correctness must be established by a reconstruction or target-factorization
identity for the retained observations. An optimizer result cannot substitute
for that identity, even when both interventions were motivated by the same
feature-merging failure.

\subsection{The Gaussian ReLU kernel from polar coordinates}
\label{qdo:kernel}

Let $(U,V)$ be jointly standard normal with correlation $\rho\in[-1,1]$.
Set $\vartheta=\arccos\rho$ and express
\begin{equation}
 U=Z_1,\qquad V=\cos\vartheta\,Z_1+\sin\vartheta\,Z_2,
 \label{qdo:normalpair}
\end{equation}
where $Z_1,Z_2$ are independent standard normals. Their polar representation is
$(Z_1,Z_2)=R(\cos\varphi,\sin\varphi)$, with $\varphi$ uniform on
$[0,2\pi)$, independent of $R$, and $\E R^2=2$. Therefore
$U=R\cos\varphi$ and $V=R\cos(\varphi-\vartheta)$.
The angles where both factors are positive form an interval equivalent to
$[\vartheta-\pi/2,\pi/2]$. It follows that
\begin{align}
 \E[\ReLU(U)\ReLU(V)]
 &=\frac{1}{\pi}\int_{\vartheta-\pi/2}^{\pi/2}
       \cos\varphi\cos(\varphi-\vartheta)\,\mathrm d\varphi\\
 &=\frac{\sin\vartheta+(\pi-\vartheta)\cos\vartheta}{2\pi}.
 \label{qdo:polarint}
\end{align}
The second equality uses
$2\cos\varphi\cos(\varphi-\vartheta)=
\cos\vartheta+\cos(2\varphi-\vartheta)$.
Thus the kernel used in the projection audit is
\begin{equation}
 \boxed{\displaystyle
 \kappa(\rho)=
 \frac{\sqrt{1-\rho^2}+(\pi-\arccos\rho)\rho}{2\pi}.}
 \label{qdo:kappa}
\end{equation}
The endpoint checks are $\kappa(1)=1/2$, $\kappa(0)=1/(2\pi)$,
and $\kappa(-1)=0$. For nonzero vectors $a,b$ and $X\sim\mathcal N(0,I_d)$,
positive homogeneity yields
\begin{equation}
 \E[\ReLU(a^\top X)\ReLU(b^\top X)]
 =\norm{a}\norm{b}\,
 \kappa\!\left(\frac{a^\top b}{\norm{a}\norm{b}}\right).
 \label{qdo:nonunitkernel}
\end{equation}
A zero vector contributes zero and requires no normalized correlation.

\subsection{Teacher moments and normalization}
\label{qdo:moments}

For a standard normal $Z$, direct integration and symmetry give
\begin{equation}
 \E\ReLU(Z)=\frac{1}{\sqrt{2\pi}},\qquad
 \E\ReLU(Z)^2=\frac12,\qquad
 \Var(\ReLU(Z))=\frac12-\frac{1}{2\pi}.
 \label{qdo:normalmoments}
\end{equation}
The four preactivations in \eqref{qdo:teacher} are independent standard
normals. Since $\sum_jc_j=0$ and $\sum_jc_j^2=1$, the target $Y=f_*(X)$
satisfies
\begin{equation}
 \E Y=0,\qquad V_*:=\Var(Y)=\frac12-\frac{1}{2\pi}.
 \label{qdo:teachermoments}
\end{equation}
The historical normalized mean squared error divides by this $V_*$
\cite{qdo,projection}. This normalization belongs to that teacher; changing
the target or input distribution requires recomputing it. The kernel
calculation below does not require a zero target mean, because it subtracts
second moments with identical means.

The positive and negative coefficients remain inside the double sum in the
general risk formula. Cross terms can have either sign, even though the final
conditional variance is nonnegative. Consequently an individual summand is
not a nonnegative amount of information assigned to one feature. The
calculation evaluates the complete target, including cancellations between
its terms. This is also why replacing the target by a visually similar sum
with different coefficients can change which contrasts must survive.

\subsection{An exact fixed-projection prediction floor}
\label{qdo:floor}

Let $P=P^\top=P^2$ be an orthogonal projector and let $Y=f_*(X)$ have finite
second moment. Define
\begin{equation}
 m_P(PX)=\E[Y\mid PX],\qquad
 R(P)=\E[\Var(Y\mid PX)].
 \label{qdo:conditionalrisk}
\end{equation}
For every square-integrable readout $\widehat Y=g(PX)$,
\begin{equation}
 \E[(Y-\widehat Y)^2]
 =R(P)+\E[(m_P(PX)-\widehat Y)^2].
 \label{qdo:orthogonalrisk}
\end{equation}
Indeed, expand $Y-\widehat Y=(Y-m_P)+(m_P-\widehat Y)$ and condition the
cross term on $PX$. It vanishes because $\E[Y-m_P\mid PX]=0$.
Thus $R(P)$ is the optimal risk over \emph{all} measurable readouts of the
fixed projection. A restricted or unsuccessfully trained decoder can incur
additional error. Changing the projection changes the floor.

For $P=P_W$, the projector onto $\im W$, observing $W^\top X$ is equivalent
to observing $P_WX$: one direction follows from $W^\top P_W=W^\top$, and the
other from $P_WX=(W^\top)^+W^\top X$. The superscript $+$ denotes the
Moore--Penrose pseudoinverse. This explains why the floor diagnoses the
information available before ReLU \cite{projection}.

The equivalence holds in exact arithmetic, regardless of the coordinates
chosen within the observed subspace. An invertible change of measurement
coordinates leaves the floor unchanged. In contrast, two projections of the
same rank can have different floors because they observe different input
directions. Ordinary rank therefore cannot distinguish all the successful and
failed task representations. Numerical conditioning affects the stability of
a reconstruction and must be evaluated separately from this exact
information comparison.

\begin{proposition}[Gaussian fixed-projection formula]
\label{qdo:floorformula}
Suppose $X\sim\mathcal N(0,I_d)$ and
$Y=\sum_i c_i\ReLU(q_i^\top X)$ with unit-norm $q_i$.
Then
\begin{equation}
 \boxed{\displaystyle
 R(P)=\sum_{i,j}c_i c_j
 \bigl[\kappa(q_i^\top q_j)-\kappa(q_i^\top Pq_j)\bigr].}
 \label{qdo:exactfloor}
\end{equation}
\end{proposition}

\begin{proof}
Take independent standard Gaussian vectors $Z,E_1,E_2$ and form
\begin{equation}
 X_1=PZ+(I-P)E_1,\qquad X_2=PZ+(I-P)E_2.
 \label{qdo:copies}
\end{equation}
Each copy has covariance $P+(I-P)=I$; their cross-covariance is $P$.
Conditioned on the shared projection $PZ$, they have independent residuals
and the correct conditional distribution of $X$. Therefore
\begin{equation}
 \E[f_*(X_1)f_*(X_2)]=\E[m_P(PX)^2].
 \label{qdo:copymoment}
\end{equation}
The pair $(q_i^\top X_1,q_j^\top X_2)$ has unit marginal variances and
correlation $q_i^\top Pq_j$. Applying \eqref{qdo:kappa} term by term gives
\begin{align}
 \E[m_P(PX)^2]&=\sum_{i,j}c_i c_j\kappa(q_i^\top Pq_j),\\
 \E[Y^2]&=\sum_{i,j}c_i c_j\kappa(q_i^\top q_j).
 \label{qdo:twomoments}
\end{align}
Finally $R(P)=\E[Y^2]-\E[m_P(PX)^2]$, proving the formula.
\end{proof}

For non-unit teacher directions, let $n_i=\norm{q_i}>0$. The extension is
\begin{equation}
 R(P)=\sum_{i,j}c_i c_j n_i n_j
 \left[
 \kappa\!\left(\frac{q_i^\top q_j}{n_i n_j}\right)
 -\kappa\!\left(\frac{q_i^\top Pq_j}{n_i n_j}\right)
 \right].
 \label{qdo:nonunitfloor}
\end{equation}
Both denominators use the \emph{original} norms. The conditional copies keep
their original marginal variances. Substituting projected norms would compute
the ReLU kernel of $Pq_i,Pq_j$, a different quantity.

An independent expression for the optimal readout is also available. With
$\mu_i=q_i^\top PX$ and $\sigma_i^2=\norm{q_i}^2-\norm{Pq_i}^2$, conditional
Gaussianity gives $q_i^\top X\mid PX\sim\mathcal N(\mu_i,\sigma_i^2)$. Hence
\begin{equation}
 m_P(PX)=\sum_i c_i
 \left[
 \sigma_i\varphi\!\left(\frac{\mu_i}{\sigma_i}\right)
 +\mu_i\Phi\!\left(\frac{\mu_i}{\sigma_i}\right)
 \right],
 \label{qdo:conditionalmean}
\end{equation}
where $\varphi,\Phi$ are the standard normal density and distribution
function. A term with $\sigma_i=0$ is interpreted as
$c_i\ReLU(\mu_i)$. This follows by integrating
$(\mu_i+\sigma_i z)$ over $z>-\mu_i/\sigma_i$.

\begin{proposition}[The required span for the original teacher]
\label{qdo:spancriterion}
For the orthonormal teacher in \eqref{qdo:teacher} and standard Gaussian
inputs, $R(P)=0$ if and only if every $q_j$ belongs to $\im P$.
\end{proposition}

\begin{proof}
If all teacher directions are observed, each preactivation and hence the
target is a function of $PX$, proving sufficiency. Conversely, zero risk
means the target is almost surely constant on almost every conditional
fiber. The conditional Gaussian has full support on that fiber. Continuity
of the ReLU teacher therefore makes it constant everywhere on each such
fiber. Those observed projections form a dense set; continuity extends this
invariance to every fiber. Thus for every $v\in\ker P$,
\begin{equation}
 f_*(x+t v)=f_*(x)\qquad\text{for every }x\text{ and }t.
 \label{qdo:fiberinvariance}
\end{equation}
Orthonormality permits a nonempty open region where exactly teacher feature
$j$ is positive. There the gradient is $c_jq_j$, so
\eqref{qdo:fiberinvariance} implies $c_jq_j^\top v=0$. Every coefficient is
nonzero, hence $q_j$ is orthogonal to $\ker P$ and belongs to $\im P$.
\end{proof}

This result is specific to this teacher with nonzero coefficients and
independently variable orthonormal features. It identifies the four required
directions without demanding preservation of the other four input
directions. It also concerns the unrectified observation: preserving that
span alone does not prove that a later nonlinear representation remains
sufficient, or that training learns the necessary decoder.

\subsection{CReLU gives an exact floor; a selective hybrid gives a bound}
\label{qdo:crelu}

CReLU retains both rectified phases of a measured vector:
\begin{equation}
 H_c=\CReLU(W^\top X)
 =\bigl(\ReLU(W^\top X),\ReLU(-W^\top X)\bigr).
 \label{qdo:creluobservation}
\end{equation}
Subtraction recovers $W^\top X$ exactly. Conversely, $H_c$ is a deterministic
function of that measurement. Consequently conditioning on $H_c$, $W^\top X$,
or $P_WX$ supplies identical information, and
\begin{equation}
 \inf_g\E[(Y-g(H_c))^2]=R(P_W).
 \label{qdo:crelufloor}
\end{equation}
This is the measurement-preservation property behind the original CReLU
filter-span reconstruction proposition, without max-pooling \cite{crelu}.
It does not restore any component in $\ker W^\top$.

Now consider the selective hybrid
\begin{equation}
 H=\bigl(\ReLU(W^\top X),\CReLU(U^\top X)\bigr),
 \qquad M=[W\ U],\qquad T=P_MX.
 \label{qdo:hybrid}
\end{equation}
Here $H$ is a function of $T$, but the reverse need not hold. Set
$m_M(T)=\E[Y\mid T]$ and $R_H=\E[\Var(Y\mid H)]$.
The conditional variance decomposition, using that $H$ is determined by $T$,
gives
\begin{align}
 R_H
 &=R(P_M)+\E[\Var(m_M(T)\mid H)]\label{qdo:hybridgap}\\
 &=R(P_M)+\E[(m_M(T)-\E[Y\mid H])^2]
 \ \ge\ R(P_M).\notag
\end{align}
For completeness, write
$Y-\E[Y\mid H]=(Y-m_M)+(m_M-\E[Y\mid H])$.
The cross term again vanishes on conditioning on $T$, and the tower property
gives $\E[m_M\mid H]=\E[Y\mid H]$. This proves the exact gap formula.

Equality holds precisely when $m_M(T)$ is recoverable from $H$ almost surely.
Recovering all of $T$ is sufficient but stronger than necessary. Merely
making $\im M$ contain the teacher span proves no such equality: the
unrectified measurements can contain a distinction that the main ReLU branch
still deletes.

Accordingly, a positive augmented-span floor rules out exact target recovery
from the hybrid. A zero augmented-span floor leaves the second term of
\eqref{qdo:hybridgap} to be checked. This term measures the target-relevant
ambiguity introduced by the actual retained observation, after all signed
measurements that were available before rectification have been taken into
account. It gives a precise test of whether a proposed selective side
channel supplies enough information.

An exact Gaussian example demonstrates a strict gap. Let
$X\sim\mathcal N(0,I_2)$, $W=e_1$, $U=e_2$, and
$Y=\ReLU(-X_1)$. Then $M=I_2$ and $R(P_M)=0$. The side channel reveals $X_2$,
which is independent of $Y$. On $X_1>0$ the target is zero; on $X_1\le0$ the
main observation is always zero and the conditional target is half-normal.
Therefore
\begin{equation}
 R_H=\tfrac12\Var(-X_1\mid X_1\le0)
 =\frac12-\frac{1}{\pi}>0.
 \label{qdo:strictgap}
\end{equation}
Thus even complete augmented linear coverage can coexist with strictly
positive optimal prediction error after selective rectification.

\subsection{Scope of the diagnostic}
\label{qdo:scope}

The conditional-risk identities hold for any square-integrable target. The
closed form \eqref{qdo:exactfloor} additionally requires Gaussian inputs,
a fixed linear observation, and a ReLU teacher of the stated form. Arbitrary
hidden states, nonlinear transition maps, pooling, quantization, or a different
input law do not automatically satisfy those assumptions. The algebraic CReLU
inverse also presumes that both measured phases remain available without
additional destructive processing.

Training the matrices first and then freezing them is compatible with the
fixed-projection calculation. Choosing a different measurement matrix as a
function of the current input is a different observation model: the selection
rule and any retained selector must be included in the conditioning
information. Substituting an input-dependent projector into the fixed-matrix
formula is not justified by the derivation above. Such adaptive records need
their own sufficiency argument or conditional-risk calculation.

Finally, $R_H=0$ states target sufficiency almost surely under the chosen
distribution. A guarantee for every admissible input requires a pointwise
factorization through the retained representation. Neither a successful probe
nor failure to find a collision supplies that proof. These distinctions let the
projection audit support the proposal as a precise diagnostic without turning
it into a general information-preservation theorem.
